\section{Sets, logic, and mappings}

The number systems used throughout these notes are
\begin{equation}
\mathbb{N}\subset\mathbb{Z}\subset\mathbb{R}\subset\mathbb{C},
\end{equation}
where $\mathbb{N}$, $\mathbb{Z}$, $\mathbb{R}$, and $\mathbb{C}$ denote the natural numbers, integers, real numbers, and complex numbers, respectively.

Logical notation makes mathematical statements compact and unambiguous. The symbols $\forall$ and $\exists$ mean ``for every'' and ``there exists.'' Thus $\forall x\in\mathbb{R}$ describes a statement valid for every real number $x$, while $\exists n\in\mathbb{N}$ asserts that at least one natural number $n$ exists with the stated property. The implication
\begin{equation}
I\Rightarrow II
\end{equation}
means that $I$ is sufficient for $II$, or equivalently that $II$ is necessary for $I$. The equivalence $I\Leftrightarrow II$ means that the two assertions are necessary and sufficient for one another. We use $I\equiv II$ when $II$ defines $I$, and $\square$ marks the end of a proof.

\begin{definition}[Mapping]
Let $A$ and $B$ be nonempty sets. A mapping from $A$ to $B$ is a rule $f$ that assigns to every $x\in A$ a unique element $y\in B$. We write
\begin{equation}
f:A\to B,\qquad y=f(x).
\end{equation}
The set $A$ is the domain and $B$ is the codomain.
\end{definition}

The mapping is injective when $f(x_1)=f(x_2)$ implies $x_1=x_2$; it is surjective when every $y\in B$ is equal to $f(x)$ for at least one $x\in A$. A bijection is both injective and surjective.

For $A\subseteq\mathbb{R}$ or $A\subseteq\mathbb{C}$, a mapping $f:A\to\mathbb{R}$ or $f:A\to\mathbb{C}$ is called a function. Its image is
\begin{equation}
f(A)=\{f(x)\mid x\in A\}.
\end{equation}
These distinctions between domain, codomain, and image will matter for inverse functions, eigenvalue problems, and complex functions.
